\documentclass[11pt]{article}
\usepackage[T1]{fontenc}
\usepackage{lmodern}
\usepackage{amsmath,amssymb}
\usepackage[margin=1in]{geometry}
\usepackage[hidelinks]{hyperref}
\setlength{\emergencystretch}{3em}

\newcommand{\Dom}{\operatorname{Dom}}
\newcommand{\dom}{\operatorname{dom}}

\begin{document}

\begin{center}
{\Large Literature-implied status: approximating zone diagrams}\\[0.4em]
{\large arXiv:1002.3583 and arXiv:1208.3124}\\[0.6em]
\textbf{Status: literature\_implied\_answer (partial subcase), not an original run proof}
\end{center}

\section*{Source questions}

The source paper is Eva Kopecka, Daniel Reem, and Simeon Reich,
\emph{Zone diagrams in compact subsets of uniformly convex normed spaces},
arXiv:1002.3583, published in \emph{Israel Journal of Mathematics} 188
(2012), 1--23 \cite{KRR2012}. In its concluding section the paper lists
several open directions: local connectedness of dominance regions, extending
existence beyond the compact/uniformly convex framework, uniqueness under
additional smoothness, removing the emanation-property hypothesis, and
approximation of zone diagrams. The branch recorded here is the approximation
branch:
\begin{quote}
``Another interesting and natural problem is how to approximate a zone diagram
in the setting described in this paper.'' \cite[Sec.~6]{KRR2012}
\end{quote}
The same paragraph notes that the proof can approximate compact sites by finite
subsets but gives no error estimates.

\section*{Supporting literature}

The supporting paper is Daniel Reem, \emph{On the Computation of Zone and
Double Zone Diagrams}, arXiv:1208.3124, published in \emph{Discrete \&
Computational Geometry} 59 (2018), 253--292 \cite{Reem2018}. It cites the
source paper as the Israel Journal of Mathematics article and treats the
problem of computing zone diagrams and double zone diagrams.

The later paper defines the inner and outer approximation sequences
$(I^{(n)})$ and $(O^{(n)})$ by iterating the $\Dom$ mapping
\cite[Algorithm~3.2]{Reem2018}. Its main convergence theorem proves that, in a
proper geodesic metric space with the geodesic inclusion property and positively
separated closed sites, the limits $m$ and $M$ are the least and greatest double
zone diagrams; in the two-site case, mixed pairs give zone diagrams
\cite[Theorem~4.3]{Reem2018}. In compact worlds, the inner and outer
approximations converge to these limits in Hausdorff distance
\cite[Corollary~4.4]{Reem2018}. The paper also explains, in the normed-space
case, how to approximate each iteration using a ray-shooting algorithm for
Voronoi/dominance regions, gives practical implementation details, and records
a time-complexity estimate \cite[Sec.~6]{Reem2018}.

This answers a concrete subcase of the 2010 approximation question: for compact
worlds inside finite-dimensional strictly convex normed spaces, and more
generally compact proper geodesic spaces with the geodesic inclusion property,
the inner/outer iterative procedure has a proved Hausdorff convergence target,
and the normed-space iterations have an explicit approximate-computation model.
When the least and greatest double zone diagrams coincide, the limit is the
unique zone diagram; the later paper notes this in particular for finite
dimensional Euclidean spaces \cite[Corollary~4.5]{Reem2018}.

\section*{Scope limitations}

This is not a full answer to every open direction in the source paper. It is a
partial literature-implied answer to the approximation branch. The later paper
does not give convergence-rate or iteration-count estimates: it explicitly says
that estimating $n_0(\epsilon)$ remains a major open problem, and its conclusion
again lists convergence-rate error estimates as open. It also does not settle
the source paper's local-connectedness, all-normed-spaces/no-compactness,
infinite-dimensional uniqueness, or no-emanation-property questions.

\section*{Search and evidence}

The cheap run indexes had no prior row for arXiv:1002.3583, zone diagrams, or
arXiv:1208.3124. The target source lines inspected were source.tex lines
765--809 of arXiv:1002.3583. The supporting source lines inspected were
arXiv:1208.3124 source lines 220--224, 362--368, 437--463, 563--665, and
1234--1247. A bounded web check confirmed the publication metadata and DOI
records for the source article and the supporting computation article.

\begin{thebibliography}{9}

\bibitem{KRR2012}
E. Kopecka, D. Reem, and S. Reich,
\emph{Zone diagrams in compact subsets of uniformly convex normed spaces},
arXiv:1002.3583; Israel J. Math. 188 (2012), 1--23.
doi:10.1007/s11856-011-0094-5.
\url{https://arxiv.org/abs/1002.3583}

\bibitem{Reem2018}
D. Reem,
\emph{On the Computation of Zone and Double Zone Diagrams},
arXiv:1208.3124; Discrete Comput. Geom. 59 (2018), 253--292.
doi:10.1007/s00454-017-9958-8.
\url{https://arxiv.org/abs/1208.3124}

\end{thebibliography}

\end{document}
